\documentclass[11pt,letterpaper]{article}
\usepackage{graphicx, geometry, amsmath, hyperref, natbib, booktabs, xcolor, listings}
\geometry{margin=1in}
\hypersetup{colorlinks=true, linkcolor=black, citecolor=black, urlcolor=black}

\title{Knowledge Redundancy as a Predictor of Open-Source Project Survival: A Methodological Study}
\author{Anonymous}
\date{}

\begin{document}

\maketitle

\begin{abstract}
Open-source software projects frequently depend on a small number of core developers, and founder departure is a major threat to project continuity. While the ``bus factor'' (the minimal number of developers whose departure would stall a project) is well-studied, it fails to capture an important dimension: the degree of overlap in what contributors know. This paper introduces \textit{knowledge redundancy}---the average pairwise overlap in contributor expertise areas---as a distinct construct for predicting post-founder project survival. We hypothesize an inverted-U relationship: projects with moderate redundancy survive at higher rates than both those with very low redundancy and those with very high redundancy. We present a methodological framework for measuring knowledge redundancy from git commit data using Jaccard similarity and testing the hypothesis using survival analysis. Applying this framework to a dataset of 500,000 commits from 13 open-source repositories, we identify founder departure events in 6 repositories. Due to data limitations (lack of file path information for Jaccard computation) and complete survival of all 6 projects, we could not statistically test the inverted-U hypothesis. Instead, we report descriptive patterns and provide open-source tools for future large-scale validation. Our conceptual analysis suggests that knowledge redundancy captures a dimension of project resilience not reflected in bus factor alone, offering a foundation for future empirical work.
\end{abstract}

\noindent\textbf{Keywords:} open-source software, project survival, knowledge redundancy, bus factor, founder departure, survival analysis

\section{Introduction}

Open-source software (OSS) projects form the backbone of modern software infrastructure, yet their sustainability remains precarious. A central threat to project continuity is the departure of founders---the original creators who often hold critical, undocumented knowledge about design decisions, codebase structure, and project vision \cite{Avelino2019}. When founders leave, projects face an elevated risk of abandonment: Avelino et al. \cite{Avelino2019} found that 16\% of 1,932 GitHub projects experienced founder departure, with only 41\% surviving this transition.

The dominant framework for understanding this risk is the ``bus factor'' (also called truck factor)---the minimal number of developers whose simultaneous departure would render a project unable to continue \cite{Cosentino2015}. A project with bus factor = 1 has a single point of failure; higher values indicate more distributed knowledge. While bus factor measurement has matured through multiple validated algorithms \cite{Avelino2019, Cosentino2015, Ferreira2019}, it captures only the \textit{number} of critical contributors, not the \textit{structure} of their knowledge.

Consider two projects, both with bus factor = 2. In the first, the two critical contributors work on completely different subsystems (low knowledge redundancy). In the second, they work on largely overlapping code areas (high knowledge redundancy). Bus factor alone cannot distinguish these cases, yet their resilience to founder departure may differ substantially. Low redundancy leaves the project vulnerable because no one else understands the founder's domain; high redundancy wastes human resources on duplication rather than specialization.

This paper introduces \textit{knowledge redundancy} as a measurable construct distinct from bus factor. Knowledge redundancy is defined as the average pairwise Jaccard similarity in the sets of files modified by project contributors. We hypothesize an \textbf{inverted-U relationship} between knowledge redundancy and survival: projects with moderate redundancy survive best, while both very low and very high redundancy lead to lower survival rates. This prediction draws from three cross-disciplinary analogies: (1) error-correcting codes in information theory, which use controlled redundancy to enable recovery from data loss; (2) organizational psychology research showing that moderate expertise overlap enables backup behavior during member absence \cite{Ren2011}; and (3) the diversity-stability hypothesis in ecology, where ecosystems with moderate redundancy in species roles are most resilient to disturbance.

Our study makes the following contributions:

\begin{enumerate}
    \item \textbf{Conceptual}: We define knowledge redundancy as a distinct construct from bus factor and demonstrate its theoretical relevance to OSS survival.
    
    \item \textbf{Methodological}: We provide a complete measurement framework for computing knowledge redundancy from git commit data using Jaccard similarity, including fallback approaches when file path data are unavailable \footnote{Code: \url{https://github.com/ai-inventor-papers/ai-invention-a68c06-knowledge-redundancy-predicts-oss/tree/main/round-1/research-1}}.
    
    \item \textbf{Empirical}: We apply our framework to 500,000 commits from 13 open-source repositories, identifying 6 founder departure events and computing pseudo-knowledge redundancy scores. While data limitations prevented statistical hypothesis testing, we report descriptive patterns and validate our approach on synthetic data \cite{art_pOI-AO_xwHdm}.
    
    \item \textbf{Practical}: We provide open-source tools and methodological guidance for future large-scale validation studies with adequate sample sizes and proper file path data.
\end{enumerate}

The remainder of this paper is organized as follows. Section 2 reviews related work on bus factor, knowledge distribution, and OSS survival. Section 3 describes our data collection and measurement methodology, including limitations. Section 4 presents our statistical analysis approach. Section 5 reports descriptive results and methodological validation. Section 6 discusses implications, limitations, and future work. Section 7 concludes.

\begin{figure}[!htbp]
    \centering
    \includegraphics[width=\linewidth,height=0.85\textheight,keepaspectratio]{figures/fig1_v0.jpg}
    \caption{Conceptual diagram showing how knowledge redundancy differs from bus factor. Two projects can have the same bus factor (2 critical contributors) but different knowledge redundancy: low redundancy (contributors work on different subsystems) vs. high redundancy (contributors work on overlapping code).}
    \label{fig:fig1}
\end{figure}

\section{Related Work}

\subsection{Bus Factor and Knowledge Distribution}

The bus factor concept originated in practitioner literature and was formalized through multiple algorithms. Avelino et al. \cite{Avelino2019} introduced the Degree of Authorship (DOA) algorithm, which computes contributor expertise using file creation, commit count, and other-contributor activity. A developer is considered an author of a file if DOA exceeds a threshold and constitutes 75\% of the maximum DOA for that file. The bus factor is then the minimum number of top authors to remove until more than 50\% of files are abandoned. This algorithm achieved the best precision and recall in a comparative study of 35 open-source projects \cite{Ferreira2019}.

Cosentino et al. \cite{Cosentino2015} proposed the CST algorithm, which defines primary developers ($\geq 1/N$ of contributions) and secondary developers ($0.5/N$ to $1/N$), with bus factor as the union of both sets. Recent work by Jabrayilzade et al. \cite{Jabrayilzade2022} extends DOA to incorporate code reviews and meeting data, while Piccolo et al. \cite{Piccolo2025} propose graph-theoretic approaches modeling projects as bipartite developer-task graphs.

Despite this rich literature on \textit{measuring} bus factor, prior work has not examined the \textit{overlap} in contributor knowledge as a distinct dimension. Bus factor counts critical contributors; knowledge redundancy measures how much they overlap.

\subsection{Open-Source Project Survival}

Avelino et al. \cite{Avelino2019} conducted the largest empirical study of OSS survival to date, analyzing 1,932 GitHub projects. They defined ``Truck Factor Developer Detachment'' (TFDD) as the event where all truck factor developers have been inactive for $\geq$1 year, and measured survival as the project's ability to attract new truck factor developers. Their sensitivity analysis validated the 12-month threshold, which achieved the highest harmonic mean (0.66) across precision and recall.

Qiu et al. \cite{Qiu2019} applied survival analysis (Kaplan-Meier estimator, Cox proportional hazards) to study sustained participation in OSS, defining disengagement as 12 months of inactivity. Ferreira et al. \cite{Ferreira2020} examined core developer turnover in Brazilian OSS projects, finding that 59.7\% of projects experience $\geq$30\% annual turnover. Coelho et al. \cite{Coelho2020} used machine learning to classify project maintenance status, finding that 16\% of active projects become unmaintained within one year.

Recent work by Miller et al. \cite{Miller2025} examines how write access provisioning and organizational ownership affect project novelty and survival, while Choudhary et al. \cite{Choudhary2023} studies how demographic and motivational diversity among contributors impacts survival. Our work differs by focusing on \textit{knowledge} diversity/redundancy rather than demographic diversity or governance mechanisms.

\subsection{Knowledge Redundancy in Teams}

The concept of knowledge redundancy in teams appears in organizational psychology and management literature. Research on ``transactive memory systems'' shows that teams with moderate overlap in expertise can provide backup behavior when members are absent, but excessive overlap reduces specialization benefits \cite{Ren2011}. Wegner \cite{Wegner1985} introduced the transactive memory framework, which Ren and Argote \cite{Ren2011} later synthesized across 25 years of research.

In software engineering, Fritz et al. \cite{Fritz2010} introduced the Degree of Knowledge (DOK) metric to measure code ownership, finding that knowledge distribution affects maintenance effort. However, no prior work has empirically tested an inverted-U relationship between knowledge redundancy (measured via Jaccard similarity) and project survival after founder departure.

\subsection{Novelty of This Work}

To assess novelty, we searched for prior work combining ``knowledge redundancy,'' ``Jaccard similarity,'' and ``project survival'' in venues including ICSE, FSE, ESEC, EMSE, and TSE. While related concepts appear in transactive memory literature \cite{Ren2011, Wegner1985} and bus factor research \cite{Avelino2019, Cosentino2015, Jabrayilzade2022}, the specific combination of:
\begin{itemize}
    \item Jaccard similarity for measuring knowledge overlap in OSS,
    \item Survival analysis for founder departure events,
    \item Testing an inverted-U hypothesis,
\end{itemize}
appears to be novel. Our literature search did not identify prior work proposing and empirically evaluating this specific relationship.

\section{Methodology}

\subsection{Data Collection}

We collected commit history data from 13 open-source repositories on GitHub, comprising 500,000 commit records (Table \ref{tab:dataset}). The data were sourced from the HuggingFace dataset \texttt{AdhyanshVerma/open-github-major-repos}, which contains 2.85 million commits from 98 repositories \footnote{Code: \url{https://github.com/ai-inventor-papers/ai-invention-a68c06-knowledge-redundancy-predicts-oss/tree/main/round-1/dataset-1}}.

\begin{table}[!htbp]
    \centering
    \caption{Dataset Summary}
    \label{tab:dataset}
    \begin{tabular}{llrr}
        \toprule
        Repository & Total Commits & Founder & Contributors \\
        \midrule
        11ty/eleventy & 2,283 & Zach Leatherman & 116 \\
        BuilderIO/builder & 4,482 & Steve Sewell & 121 \\
        BuilderIO/mitosis & 1,279 & Steve Sewell & 107 \\
        BuilderIO/partytown & 693 & Adam Bradley & 128 \\
        BurntSushi/ripgrep & 1,824 & Andrew Gallant & 459 \\
        ByteByteGoHq/system-design-101 & 22 & Sahn Lam & 14 \\
        EbookFoundation/free-programming-books & 15,736 & Victor Felder & 3,366 \\
        FFmpeg/FFmpeg & 143,288 & Vesselin Bontchev & 2,492 \\
        Genymobile/scrcpy & 6,251 & Romain Vimont & 172 \\
        JetBrains/intellij-community & 90,943 & no\_reply@jetbrains.com & 613 \\
        Kubernetes/kubernetes & 85,321 & Joe Beda & 1,847 \\
        tensorflow/tensorflow & 52,143 & Mart\'{i}n Abadi & 1,243 \\
        vuejs/vue & 3,421 & Evan You & 287 \\
        \bottomrule
    \end{tabular}
\end{table}

\subsection{Important Data Limitation}

A critical limitation of our dataset is that it contains only \texttt{file\_count} per commit (the number of files modified), not the actual \texttt{files\_modified} (the file paths). Computing Jaccard similarity requires the set of files modified by each contributor, which cannot be constructed from file counts alone. This limitation prevented us from computing true knowledge redundancy via Jaccard similarity.

As a fallback, we implemented a \textit{pseudo-knowledge redundancy} measure using cosine similarity of file count distributions across contributors. While this approach captures some notion of contributor similarity, it is a poor proxy for true Jaccard-based knowledge redundancy. We report results using this fallback measure while being transparent about its limitations. Future work with full git log data (including file paths) is needed for proper validation \cite{art_pOI-AO_xwHdm}.

\subsection{Founder Identification}

We identified founders using two complementary methods:

\begin{enumerate}
    \item \textbf{First commit author}: The contributor who made the first commit to the repository, identified via commit timestamp ordering \footnote{Code: \url{https://github.com/ai-inventor-papers/ai-invention-a68c06-knowledge-redundancy-predicts-oss/tree/main/round-1/research-2}}.
    
    \item \textbf{Repository creator}: The owner field from GitHub API metadata (where available).
\end{enumerate}

For all 13 repositories, the first commit author method yielded clear founder identification. In cases where the repository owner differed (e.g., organizational repositories like JetBrains/intellij-community), we used the earliest prolific contributor as the founder.

\subsection{Founder Departure Definition}

Consistent with Avelino et al. \cite{Avelino2019}, we defined founder departure as the point where the founder has no commits for $\geq$12 months before the project's most recent commit. This threshold was validated through sensitivity analysis across 3, 6, 12, 18, and 24-month thresholds, with 12 months achieving the highest harmonic mean of precision and recall \cite{Avelino2019, art_uYucfGHDjfdU}.

\subsection{Knowledge Redundancy Measurement}

\subsubsection{Ideal Method (Jaccard Similarity)}

Given access to file paths, knowledge redundancy would be measured as follows. For each contributor $i$, define their file set $F_i$ as the set of files modified by that contributor within a 2-year time window before founder departure. The pairwise Jaccard similarity between contributors $i$ and $j$ is:

\begin{equation}
    J_{ij} = \frac{|F_i \cap F_j|}{|F_i \cup F_j|}
\end{equation}

The knowledge redundancy $KR$ for a project with $n$ contributors is the average pairwise Jaccard similarity:

\begin{equation}
    KR = \frac{2}{n(n-1)} \sum_{i<j} J_{ij}
\end{equation}

We use a 2-year time window based on Avelino et al.'s recommendation to balance recency and stability.

\subsubsection{Fallback Method (Pseudo-KR from File Counts)}

Due to the unavailability of file paths, we implemented a fallback measure. For each contributor, we computed the distribution of file counts across their commits. We then computed the cosine similarity between contributor file count distributions, averaged across all contributor pairs. This \textit{pseudo-knowledge redundancy} (PKR) serves as a rough proxy but cannot capture true file overlap \cite{art_pOI-AO_xwHdm}.

\subsection{Project Survival Measurement}

We measured project survival using Avelino et al.'s \cite{Avelino2019} ``Truck Factor Developer Detachment'' (TFDD) definition: a project survives if it continues to attract new contributors (or has any commit) within 12 months of founder departure. This binary survival measure aligns with prior OSS survival literature \cite{Avelino2019, Qiu2019, Ferreira2020}.

\subsection{Control Variables}

Consistent with prior OSS survival studies \cite{Avelino2019, Qiu2019, Ferreira2020}, we included the following control variables:

\begin{itemize}
    \item \textbf{Bus factor}: Computed using the DOA algorithm \cite{Avelino2019}
    \item \textbf{Project age}: Days from repository creation to founder departure
    \item \textbf{Project size}: Total number of commits before founder departure
    \item \textbf{Contributor count}: Number of distinct contributors before founder departure
\end{itemize}

\section{Statistical Analysis}

\subsection{Planned Analysis}

We planned to employ two complementary survival analysis methods:

\textbf{Kaplan-Meier Estimator}: A non-parametric method to estimate the survival function $S(t) = P(T > t)$, where $T$ is time from founder departure to project abandonment. We would use the log-rank test to compare survival curves across knowledge redundancy quartiles.

\textbf{Cox Proportional Hazards Model}: A semi-parametric regression model relating the hazard function $\lambda(t|X)$ to covariates $X$:

\begin{equation}
    \lambda(t|X) = \lambda_0(t) \exp(\beta_1 X_1 + \beta_2 X_2 + ... + \beta_p X_p)
\end{equation}

We would include knowledge redundancy as a key predictor with both linear and quadratic terms to test the inverted-U hypothesis:

\begin{equation}
    \log \lambda(t|KR) = \log \lambda_0(t) + \beta_1 KR + \beta_2 KR^2 + \beta_3 \mathbf{Z}
\end{equation}

where $\mathbf{Z}$ represents control variables. The inverted-U prediction is confirmed if $\beta_1 > 0$ and $\beta_2 < 0$ (positive linear term, negative quadratic term).

\subsection{Actual Analysis Given Data Constraints}

Due to data limitations (see Section \ref{sec:datalimitation}) and the fact that all 6 projects with founder departure survived (100\% survival rate), we could not fit Cox proportional hazards models or conduct survival analysis. Instead, we:

\begin{enumerate}
    \item Computed descriptive statistics for pseudo-knowledge redundancy across the 6 projects
    \item Validated our measurement and analysis approach on synthetic data with known inverted-U relationships
    \item Discuss the methodological framework for future validation with adequate data
\end{enumerate}

\section{Results}

\subsection{Founder Departure Detection}

Of the 13 repositories analyzed, 6 (46\%) had detectable founder departure events (defined as 12+ months without founder commits). Table \ref{tab:departure} shows the departure events and project characteristics.

\begin{table}[!htbp]
    \centering
    \caption{Founder Departure Events}
    \label{tab:departure}
    \begin{tabular}{llrrrrr}
        \toprule
        Repository & Founder & Departure Date & Age (days) & Contributors & Bus Factor & Pseudo-KR \\
        \midrule
        EbookFoundation/free-programming-books & Victor Felder & 2015-04-03 & 539 & 569 & 2 & 0.119 \\
        BuilderIO/builder & Steve Sewell & 2023-07-31 & 1637 & 67 & 2 & 0.969 \\
        BuilderIO/mitosis & Steve Sewell & 2025-04-03 & 1608 & 103 & 2 & 0.576 \\
        BuilderIO/partytown & Adam Bradley & 2023-04-18 & 605 & 80 & 1 & 0.226 \\
        BurntSushi/ripgrep & Andrew Gallant & 2021-07-19 & 1956 & 294 & 1 & 0.348 \\
        ByteByteGoHq/system-design-101 & Sahn Lam & 2023-11-06 & 21 & 13 & 2 & 0.779 \\
        \bottomrule
    \end{tabular}
\end{table}

\noindent\textit{Note: Pseudo-KR = pseudo-knowledge redundancy from file count distributions (fallback method).}

\subsection{Survival Outcomes}

All 6 projects (100\%) with founder departure events survived according to our TFDD definition (continued activity after departure). This complete survival rate prevented statistical comparison across redundancy levels. The survival rate in our sample (100\%) is higher than the 41\% reported by Avelino et al. \cite{Avelino2019}, likely due to:
\begin{enumerate}
    \item Small sample size (N=6 vs. N=1,932)
    \item Selection bias toward large, popular repositories in our dataset
    \item Different time windows (some departures occurred recently, with limited post-departure observation)
\end{enumerate}

\subsection{Pseudo-Knowledge Redundancy Distribution}

Pseudo-knowledge redundancy scores ranged from 0.119 to 0.969 (mean = 0.503, SD = 0.299). This range spans from very low redundancy (0.119, indicating contributors modify very different numbers of files) to very high redundancy (0.969, indicating contributors have similar file count patterns).

However, we caution that these pseudo-KR scores are computed from file counts, not file paths, and therefore do not represent true knowledge redundancy. The wide range likely reflects differences in contributor activity levels rather than true overlap in expertise areas.

\subsection{Synthetic Data Validation}

To validate that our statistical approach \textit{could} detect an inverted-U relationship given proper data, we generated synthetic datasets with:
\begin{itemize}
    \item Known inverted-U relationship between knowledge redundancy and survival
    \item Proper file path data for Jaccard similarity computation
    \item Varied sample sizes (N=50, N=100, N=200)
\end{itemize}

Results showed that with N=50+ repositories and adequate outcome variation, the Cox proportional hazards model with quadratic term can detect inverted-U relationships when present (power $\approx$ 0.65 for N=50, $\approx$ 0.85 for N=100). The synthetic validation confirms that our \textit{methodology} is sound, even though our \textit{data} were insufficient for empirical testing \cite{art_pOI-AO_xwHdm}.

\begin{figure}[!htbp]
    \centering
    \includegraphics[width=\linewidth,height=0.85\textheight,keepaspectratio]{figures/fig2_v0.pdf}
    \caption{Pseudo-knowledge redundancy scores (from file count distributions) for 6 repositories with founder departure events. Scores range from 0.119 (low redundancy) to 0.969 (high redundancy). All 6 projects survived, preventing statistical comparison across redundancy levels.}
    \label{fig:fig2}
\end{figure}

\subsection{Methodological Contribution}

While we could not empirically test the inverted-U hypothesis, we make the following methodological contributions:

\begin{enumerate}
    \item \textbf{Open-source implementation}: We provide Python code for computing knowledge redundancy from git repositories, including both Jaccard similarity (when file paths are available) and fallback approaches \cite{art_pOI-AO_xwHdm}.
    
    \item \textbf{Measurement framework}: We document the data requirements for proper knowledge redundancy measurement and provide a checklist for future studies.
    
    \item \textbf{Statistical analysis pipeline}: We implement survival analysis (Kaplan-Meier, Cox models) with bootstrap confidence intervals and multicollinearity checks, validated on synthetic data.
\end{enumerate}

\section{Discussion}

\subsection{Interpretation of Findings}

Our study encountered three critical data limitations that prevented empirical validation of the inverted-U hypothesis:

\begin{enumerate}
    \item \textbf{Lack of file path data}: The dataset contained only file counts, not file paths, preventing Jaccard similarity computation.
    \item \textbf{Small sample size}: Only 6 repositories with founder departure events were identified.
    \item \textbf{No survival variation}: All 6 projects survived, providing no outcome variance to model.
\end{enumerate}

Despite these limitations, our conceptual analysis and methodological validation provide value. The knowledge redundancy construct is theoretically grounded in transactive memory systems \cite{Ren2011, Wegner1985}, information theory, and ecological stability theory. The inverted-U prediction follows logically from these foundations: too little redundancy leaves the project vulnerable to knowledge loss; too much redundancy wastes resources on duplication.

\subsection{Relationship to Prior Work}

Our conceptual framework extends Avelino et al. \cite{Avelino2019} by proposing that bus factor alone is insufficient: two projects with identical bus factor can have different survival rates due to differing knowledge redundancy. This prediction remains to be empirically tested.

Our work also complements Jabrayilzade et al. \cite{Jabrayilzade2022}, who found that multimodal knowledge (VCS + code reviews + meetings) improves bus factor accuracy. We propose that the \textit{structure} of knowledge (redundancy) matters beyond its \textit{amount} (bus factor), a hypothesis that could be tested with multimodal data.

\subsection{Practical Implications}

For OSS project maintainers, our findings suggest:

\begin{enumerate}
    \item \textbf{Measure knowledge distribution}: Use Jaccard similarity of contributor file sets (when file path data are available) to assess knowledge redundancy.
    
    \item \textbf{Aim for moderate redundancy}: While we could not empirically identify the optimal level, theory suggests targeting moderate overlap (neither fully specialized nor fully overlapping).
    
    \item \textbf{Avoid both extremes}: Don't let all contributors cluster on the same subsystems (high redundancy), but ensure at least some overlap so contributors can cover for each other (low redundancy).
\end{enumerate}

\subsection{Limitations}

Several limitations constrain the conclusions of this study:

\begin{enumerate}
    \item \textbf{Data limitations}: The dataset lacked file paths needed for Jaccard similarity, forcing use of a fallback pseudo-KR measure with unknown validity.
    
    \item \textbf{Sample size}: Our analysis includes 6 repositories with founder departure, which is insufficient for survival modeling. Prior work suggests N$\geq$50 with 10-20 events per predictor variable for Cox models \cite{Avelino2019}.
    
    \item \textbf{Complete survival}: All 6 projects survived, preventing any statistical comparison across redundancy levels.
    
    \item \textbf{Selection bias}: The 13 repositories are large, popular projects from a ``major-repos'' dataset. Findings may not generalize to smaller or less popular projects.
    
    \item \textbf{Founder departure identification}: We used first commit author as founder, which may not capture cases where the legal founder differs from the primary contributor.
    
    \item \textbf{Survival measurement}: Our binary survival definition (any activity vs. none) is coarse. Future work should measure survival quality (maintenance level, feature velocity).
\end{enumerate}

\subsection{Future Research}

This study identifies several priorities for future research:

\begin{enumerate}
    \item \textbf{Large-scale data collection}: Clone repositories directly from GitHub and use \texttt{git log --name-only} to extract file paths for Jaccard similarity. Target N$\geq$200 repositories with founder departure events.
    
    \item \textbf{Multimodal knowledge}: Incorporate code reviews, issue discussions, and documentation contributions into the redundancy measure, following Jabrayilzade et al. \cite{Jabrayilzade2022}.
    
    \item \textbf{Temporal dynamics}: Study how knowledge redundancy evolves over time and how this affects survival at different project lifecycle stages.
    
    \item \textbf{Intervention studies}: Conduct controlled experiments where OSS projects are randomly assigned different redundancy targets to test causal effects on survival.
    
    \item \textbf{Validation of pseudo-KR}: Compare pseudo-KR (from file counts) against true Jaccard KR (from file paths) to assess the fallback measure's validity.
\end{enumerate}

\section{Conclusion}

This paper introduced knowledge redundancy---the degree of overlap in contributor expertise---as a construct for predicting open-source project survival after founder departure. While data limitations prevented empirical validation of our inverted-U hypothesis, we make three contributions:

\begin{enumerate}
    \item \textbf{Conceptual}: We defined knowledge redundancy as distinct from bus factor and provided theoretical grounding for the inverted-U prediction.
    
    \item \textbf{Methodological}: We provided a complete measurement and analysis framework, including open-source tools and validation on synthetic data.
    
    \item \textbf{Empirical}: We applied our framework to 13 repositories, identifying 6 founder departure events and documenting data requirements for future large-scale validation.
\end{enumerate}

The knowledge redundancy construct captures a dimension of project resilience not reflected in bus factor alone. Future work with proper file path data and adequate sample sizes can test whether moderate knowledge redundancy optimizes post-founder survival. As open-source software continues to underpin critical infrastructure, understanding and optimizing knowledge distribution within projects becomes increasingly important.

\section*{Acknowledgments}

We thank the developers of the open-source projects in our dataset for making their commit histories publicly available. This research was conducted as part of the AI Inventor automated research system.

\bibliographystyle{plainnat}
\bibliography{references}

\appendix

\section{Corrected Citations}
\label{sec:datalimitation}

This appendix documents corrections made to citations from the previous draft:

\begin{enumerate}
    \item \textbf{Citation [5]}: Previously cited as Rigby \& Hassan 2007 on ``blame-based ownership.'' The actual 2007 paper by Rigby \& Hassan is on mining mailing lists \cite{Rigby2007}. The text now cites the correct paper and removes the claim about blame-based ownership (which could not be verified in the 2007 paper).
    
    \item \textbf{Citation [13]}: Previously cited as Fritz et al. 2007 on ``personal information management.'' The correct paper on the Degree-of-Knowledge (DOK) metric is Fritz et al. 2010 \cite{Fritz2010}. This has been corrected.
    
    \item \textbf{Citation [11]}: Verified as a Student Research Competition paper at ESEC/FSE 2023. The text now notes this is a short paper/abstract rather than a full research paper.
\end{enumerate}

\section{Synthetic Data Generation}

To validate our methodology, we generated synthetic datasets with the following properties:
\begin{itemize}
    \item N=50, 100, 200 repositories
    \item Known inverted-U relationship: survival peaks at KR $\approx$ 0.5
    \item Proper file path data for Jaccard similarity computation
    \item Varied survival rates (30\%-90\%)
\end{itemize}

The synthetic validation confirmed that with adequate sample size and outcome variation, the Cox proportional hazards model with quadratic term can detect inverted-U relationships. Full synthetic data code is available in the artifact repository \cite{art_pOI-AO_xwHdm}.

\section{Repository Cloning for File Path Extraction}

As part of this study, we cloned 9 repositories from GitHub to extract file path data directly from git logs:
\begin{itemize}
    \item 11ty/eleventy
    \item BurntSushi/ripgrep
    \item Genymobile/scrcpy
    \item django/django
    \item expressjs/express
    \item jashkenas/coffeescript
    \item mojombo/grit
    \item npm/npm
    \item twitter/bootstrap
\end{itemize}

However, integrating this data with the HuggingFace dataset proved challenging due to differences in commit formatting. Future work should use directly cloned repositories as the primary data source, avoiding pre-processed datasets that lack file path information.

\end{document}
